\begin{textsample}{POS Dim 1 – vem\_ed\_20 – Score 8.00 – vem\_ed\_20/t101.txt}  \label{ex:f1_pos_110}
“Para a longevidade dos projetos, \textbf{é} importante ter a crença nos benefícios \textbf{que} trazem para professores, alunos e instituição; mostrar-se confiável ao parceiro por meio de uma comunicação intensa e honesta; ter \textbf{flexibilidade}, \textbf{compreensão} das características culturais do colega estrangeiro e otimismo”, resumiu Moreira.
Na última década, os PCIs envolveram 50 Fatecs, 7.484 alunos e 49 instituições de ensino \textbf{superior} em 17 países.
Entre 2018 e 2022, a Unesp realizou 95 projetos em 27 unidades com 54 instituições parceiras em 17 países, envolvendo 1.924 alunos. “Para garantir a longevidade do \textbf{programa} de Intercâmbios Virtuais, \textbf{é} importante informar a comunidade acadêmica, oferecer formação de docentes, auxiliar na busca de parceiros internacionais, \textbf{reconhecer} as \textbf{ações} de internacionalização no âmbito \textbf{institucional} e registrar as atividades de Intercâmbio Virtual nos \textbf{históricos} escolares dos alunos”. “Intercâmbios Virtuais e \textbf{extensão} com a comunidade” \textbf{foram} abordados pelos professores Jan Krimphove (coordenador de assuntos internacionais da Unichristus / Ceará) e Silvio Ribeiro (professor da Fatec Lins).
A universidade cearense desenvolve colaborações com a DePaul University (EUA) na área de saúde. “Se bem planejados, os Intercâmbios Virtuais podem contribuir para as atividades de \textbf{extensão}; além disso, \textbf{são} importantes oportunidades de internacionalização para desenvolver a cidadania global dos estudantes”, comentou Jan Krimphove.

% matched lemmas: ação, compreensão, extensão, flexibilidade, histórico, institucional, programa, que, reconhecer, ser, superior
\end{textsample}
